\chapter[Proposta do Protótipo]{Proposta do Protótipo}
\label{cap:prototipo}

Este capítulo descreve a proposta de solução para a extração de relações clínicas.
O trabalho integra o desenvolvimento do \textbf{RECLin-PT}, um modelo aprimorado de
extração de relações em texto clínico em português. Como etapa inicial dessa
construção, conduz-se um \textbf{experimento controlado de \emph{baselines}} que
responde a uma pergunta de projeto anterior a qualquer aprimoramento: \emph{o
pré-treinamento clínico importa para a extração de relações em textos médicos em
português?} A resposta determina sobre qual \emph{encoder} o RECLin-PT deve ser
construído. Respondida essa pergunta, a Seção~\ref{sec:reclin} descreve a segunda
etapa, em que o RECLin-PT é construído sobre a base assim escolhida por meio de um
filtro de pista lexical calibrado no conjunto de desenvolvimento.

\section{Visão geral e papel dos \emph{baselines}}
\label{sec:visao_baselines}

Comparam-se dois \emph{encoders} para o português, tratados como \emph{baselines}
sob a mesma tarefa: o \textbf{BioBERTpt} \cite{schneider_2020_biobertpt}, adaptado
a narrativas clínicas e biomédicas em português, e o \textbf{BERTimbau}
\cite{souza_2020_bertimbau}, pré-treinado em texto de domínio geral (brWaC). Os
dois compartilham a mesma arquitetura base (BERT) e diferem essencialmente no
corpus de pré-treinamento, o que os torna candidatos naturais a um experimento
controlado sobre a importância do domínio de pré-treino.

A premissa metodológica central é a \textbf{paridade}. A infraestrutura descrita
aqui não constitui um modelo com nome próprio, e sim a cadeia que viabiliza a
comparação de forma justa. Os dois \emph{baselines} compartilham um único núcleo de
implementação (\texttt{src/relation\_extraction.py}) e são, por construção,
idênticos em tudo exceto no \emph{checkpoint} carregado:

\begin{itemize}
  \item a mesma partição congelada e os mesmos pares candidatos
        (Capítulo~\ref{cap:metodologia}), o que coloca os dois sob o mesmo teto de
        \emph{recall} e na mesma régua estatística;
  \item a mesma representação de entrada por marcadores de entidade
        (Seção~\ref{sec:markers});
  \item a mesma função de perda, \emph{cross-entropy} com
        \texttt{class\_weight="balanced"};
  \item o mesmo escalonador linear com \emph{warmup} e os mesmos
        hiperparâmetros de otimização (Quadro~\ref{quad:ambiente});
  \item a mesma semente de treino e o mesmo critério de seleção de modelo, o maior
        macro-F1 no conjunto de validação.
\end{itemize}

Trocar um \emph{baseline} pelo outro consiste, literalmente, em mudar a
\emph{string} do \emph{checkpoint}. Qualquer diferença observada pode, portanto,
ser atribuída ao pré-treinamento, e não a artefatos de implementação. Essa é
exatamente a condição para que a pergunta de pesquisa tenha resposta.

\section{Os dois \emph{encoders} comparados}

\subsection{BioBERTpt --- \emph{encoder} de domínio clínico}

O \textbf{BioBERTpt} (\texttt{pucpr/biobertpt-all})\footnote{\emph{Checkpoint} público
em \href{https://huggingface.co/pucpr/biobertpt-all}{huggingface.co/pucpr/biobertpt-all}.}
\cite{schneider_2020_biobertpt}
é um modelo contextual adaptado ao domínio clínico e biomédico em português
brasileiro, com $177{,}9$ milhões de parâmetros. Sua adoção corresponde à
expectativa, difundida na literatura e sustentada pelos ganhos relatados em
reconhecimento de entidades nomeadas clínicas, de que o pré-treinamento sobre
texto do próprio domínio produza representações mais adequadas a tarefas clínicas.
É essa expectativa que o experimento coloca à prova.

\subsection{BERTimbau --- \emph{encoder} de domínio geral}

O \textbf{BERTimbau}
(\texttt{neuralmind/bert-base-portuguese-cased})\footnote{\emph{Checkpoint} público em
\href{https://huggingface.co/neuralmind/bert-base-portuguese-cased}{huggingface.co/neuralmind/bert-base-portuguese-cased}.}
\cite{souza_2020_bertimbau} é a referência de domínio geral para o português,
pré-treinado no corpus brWaC, com $108{,}9$ milhões de parâmetros, cerca de
$60\%$ do tamanho do BioBERTpt. Ele entra no experimento como o controle que isola
o efeito do domínio de pré-treino, e não como modelo secundário. Se o desempenho
for equivalente, o modelo menor é preferível por custo de inferência e de memória,
consideração concreta para implantação em ambientes clínicos com recursos
limitados.

\section{Tratamento do desbalanceamento}

Uma decisão central, comum aos dois \emph{baselines}, é o uso de
\texttt{class\_weight="balanced"} na função de perda, que reescala o erro de forma
inversamente proporcional à frequência de cada classe. Sem esse ajuste, o
otimizador minimizaria o erro simplesmente prevendo \texttt{no\_relation} para
quase tudo, dado que essa classe concentra $94{,}02\%$ dos candidatos de teste. A
reescala força o modelo a ``pagar'' caro por erros nas classes raras, sem descartar
dados nem introduzir o não-determinismo de uma subamostragem. Como a mesma
estratégia é aplicada aos dois \emph{encoders}, seu efeito é compartilhado e não
confunde a comparação.

A seleção de modelo usa o conjunto de validação (macro-F1), com o conjunto de teste
reservado a uma única avaliação final. Os hiperparâmetros efetivamente utilizados
constam do Apêndice~\ref{ap:hiperparametros}.

\section{Representação de entrada: marcadores de posição}
\label{sec:markers}

O modelo contextual não recebe o par de entidades de forma abstrata, e sim um
recorte do \emph{texto} do documento com o par delimitado por \textbf{marcadores
de entidade}. A técnica é padrão em extração de relações com modelos baseados em
\foreignlanguage{english}{BERT} \cite{soares_2019_markers}. Inserir no próprio texto
marcadores que delimitam $e_1$ e $e_2$ é mais eficaz do que concatenar
representações de \emph{span} isoladas, pois permite que o mecanismo de atenção
condicione a codificação de toda a janela ao par sob análise.

\subsection{O esquema adotado}

A função \texttt{build\_marked\_window} (\texttt{src/relation\_extraction.py})
recorta uma janela de $\pm 128$ caracteres em torno do par
(\texttt{ctx\_chars}) e insere quatro delimitadores, no esquema
%
\begin{center}
\texttt{[E1] superfície$_{e_1}$ [/E1]} \quad\dots\quad
\texttt{[E2] superfície$_{e_2}$ [/E2]}
\end{center}
%
Os quatro marcadores são registrados como \emph{tokens}\footnote{\emph{Tokens} são
as unidades em que o \emph{tokenizer} divide o texto antes de entregá-lo ao
modelo, e podem ser palavras inteiras ou fragmentos de palavra. Os especiais não
vêm do texto e recebem um vetor próprio.} especiais no
\emph{tokenizer} (\texttt{add\_special\_tokens}), e a matriz de \emph{embeddings}
do modelo é redimensionada (\texttt{resize\_token\_embeddings}) para lhes dar um
vetor dedicado, aprendido durante o ajuste fino.

Três propriedades desse esquema merecem registro explícito, por delimitarem
exatamente \emph{que informação} o modelo recebe:

\begin{enumerate}
  \item \textbf{Os marcadores são posicionais, não semânticos.} Eles indicam
        quais trechos da janela formam o par e qual deles é o primeiro
        argumento, e nada mais. O \textbf{tipo semântico} anotado pelo SemClinBr
        (\texttt{Disease or Syndrome}, \texttt{Negation}, \texttt{Finding}\ldots)
        \textbf{não é passado ao modelo}, e cabe a ele inferir a natureza de cada
        entidade a partir da superfície textual e do contexto.
  \item \textbf{A direção vem da etiqueta, não da ordem.} A inserção é feita por
        deslocamento de caractere; quando $e_2$ precede $e_1$ no texto, a
        \emph{string} resultante começa por \texttt{[E2]}. A direção da
        relação, que importa para \texttt{negation\_of}, é preservada por
        \emph{qual} marcador envolve \emph{qual} entidade, e não pela ordem em
        que os marcadores aparecem.
  \item \textbf{A classificação usa a representação agregada da sequência.} O
        cabeçote é o de \texttt{AutoModelForSequenceClassification}, que projeta
        a representação agregada da sequência (o \emph{token} \texttt{[CLS]}) nas
        três classes. Não há extração explícita dos estados ocultos dos
        marcadores. Eles influenciam a predição por meio da atenção, ao
        condicionar a codificação da janela inteira ao par delimitado.
\end{enumerate}

Como a mesma representação é aplicada aos dois \emph{encoders}, qualquer efeito
dessas escolhas é compartilhado e não confunde a comparação
(Capítulo~\ref{cap:experimentos}).

\subsection{Exemplo}

Aplicado ao trecho do Quadro~\ref{quad:exemplo_candidatos}, o par \emph{gold}
$16 \rightarrow 52$ produz a entrada reproduzida no Quadro~\ref{quad:entrada}.
Note que nem \texttt{Disease or Syndrome} (tipo de ``OSTEOMIELITE'') nem
\texttt{Negation} (tipo de ``NÃO ESPECIFICADA'') aparecem na entrada, pois o modelo
vê apenas o texto e os delimitadores.

\begin{quadro}[htbp]
  \centering
  \caption{Entrada efetivamente construída para o par \emph{gold} do Quadro~\ref{quad:exemplo_candidatos}.}
  \label{quad:entrada}
  \begin{tabular}{p{0.92\textwidth}}
  \toprule
  \texttt{\ldots{} CURATIVO COM ASPECTO LIMPO E SECO. SEGUE EM OBSERVAÇÃO.
  [E1] OSTEOMIELITE [/E1] [E2] NÃO ESPECIFICADA [/E2]. INTUBADO, SOB VENTILAÇÃO
  MECÂNICA, MONITORIZAÇÃO CARDÍACA \ldots} \\
  \bottomrule
  \end{tabular}\\[4pt]
  {\footnotesize Fonte: elaborado pelo autor (saída de \texttt{build\_marked\_window} para o documento 8906 do SemClinBr, \texttt{ctx\_chars}$=128$).}
\end{quadro}

\subsection{O que essa escolha deixa de fora}
\label{sec:markers_limitacao}

Registrar o que \emph{não} entra na representação é tão importante quanto
descrever o que entra. O SemClinBr anota cada entidade com tipos semânticos da
UMLS, frequentemente em \textbf{rótulos múltiplos} separados por ``\texttt{|}''
(por exemplo, \texttt{Medical Device|Finding|Abbreviation}). Essa anotação
\textbf{não chega ao modelo}. Os resultados do Capítulo~\ref{cap:experimentos}
foram obtidos com marcadores puramente posicionais, e a natureza de cada entidade
tem de ser inferida da superfície textual e do contexto.

A anotação de tipo, aliás, não é diretamente aproveitável na forma em que vem. No
conjunto de treino ocorrem apenas $89$ tipos \emph{atômicos} distintos, mas a
combinação multirrótulo faz o número de \emph{strings} de tipo distintas subir para
$828$, das quais $752$ contêm ``\texttt{|}''. Injetar essa \emph{string}
diretamente no marcador seria contraproducente, pois cada uma das $828$ combinações
viraria um símbolo praticamente único e mal povoado no treino, prejudicando a
generalização. Aproveitá-la exigiria, antes, um esquema de normalização que
colapsasse as combinações em um vocabulário pequeno e estável.

Para a pergunta desta etapa, porém, a ausência do tipo é inclusive conveniente.
Como nenhum dos dois \emph{encoders} recebe a informação semântica pronta, a
comparação mede a capacidade de cada um de \emph{inferi-la} do texto, e é
exatamente aí que o pré-treinamento de domínio deveria fazer diferença.

\section{Construção do RECLin-PT}
\label{sec:reclin}

Respondida a pergunta de projeto, esta seção descreve a segunda etapa da
construção. A especialização é dirigida à classe \texttt{negation\_of}, pelo
critério de criticidade clínica enunciado na Seção~\ref{sec:motivacao}, e toma a
forma de um \textbf{filtro de pista lexical} aplicado às predições do
\emph{encoder} já ajustado.

\subsection{O diagnóstico que motiva o filtro}

O ponto de partida está nas colunas de precisão e \emph{recall} da classe-alvo.
Nas quatro execuções dos \emph{baselines}, o \emph{recall} de
\texttt{negation\_of} varia entre $0{,}9145$ e $0{,}9342$, ao passo que a precisão
fica entre $0{,}5547$ e $0{,}6167$ (Tabela~\ref{tab:fase2_resultados}). O gargalo
da classe é, portanto, de \textbf{precisão}, e não de cobertura. O modelo localiza
o contexto de negação e o atribui ao par errado, produzindo falsos positivos em
candidatos que sequer começam por uma expressão de negação. Um ganho nessa classe
depende, por esse motivo, de restringir \emph{onde} ela pode ser predita, e não de
deslocar a fronteira entre \texttt{negation\_of} e \texttt{associated\_with}.

\subsection{Indução do léxico de gatilhos}

Toda relação \texttt{negation\_of} anotada no SemClinBr tem, como primeiro
argumento, uma expressão que marca a negação (``sem'', ``nega'', ``não'',
``ausência'', além das negações morfológicas registradas no
Apêndice~\ref{ap:lexico}). O léxico de gatilhos empregado pelo filtro
(\texttt{src/negation\_lexicon.py}) é induzido \textbf{exclusivamente do conjunto
de treino}. Induzi-lo do corpus completo configuraria vazamento de dados, pela
mesma razão que motiva o particionamento em nível de documento
(Seção~\ref{sec:particao}), e comprometeria a leitura de qualquer ganho medido no
conjunto de teste.

A indução percorre o \emph{espaço de candidatos} do treino, gerado por
\texttt{iter\_candidate\_pairs} sob a mesma janela \texttt{max\_gap}$=25$
(Seção~\ref{sec:candidatos}), e não o conjunto de relações anotadas. Os dois
diferem, uma vez que parte das relações \emph{gold} fica fora da janela e nunca
chega a ser candidata (Seção~\ref{sec:teto_recall}). Como o filtro opera sobre
predições feitas no espaço de candidatos, o léxico é induzido desse mesmo espaço,
de modo que ambos falem do mesmo objeto. São contadas as formas de superfície de
$e_1$ nos pares candidatos rotulados \texttt{negation\_of}, e retêm-se aquelas
cuja frequência atinge o limiar \texttt{min\_freq}.

As formas são comparadas sob normalização por decomposição NFKD, remoção de
diacríticos, colapso de espaço em branco e \emph{casefold}. A remoção de
diacríticos não é cosmética. No corpus, a mesma forma ocorre acentuada e não
acentuada (``NÃO'' e ``NAO'', ``AUSÊNCIA'' e ``AUSENCIA''), e sem essa
normalização um único gatilho se partiria em duas formas, cada uma delas abaixo
do limiar de frequência por artefato de digitação.

\subsection{A regra de rebaixamento}

A regra de decisão atua sobre predições já produzidas e é de \textbf{rebaixamento}.
Toda predição \texttt{negation\_of} cujo primeiro argumento não pertença ao léxico
é reescrita como \texttt{no\_relation}. Nenhuma predição é promovida a
\texttt{negation\_of}, e nenhuma outra classe é tocada.

Dessa assimetria decorrem duas consequências que importam para a leitura dos
resultados. A primeira é que o F1 de \texttt{associated\_with} permanece inalterado
em todos os dígitos nas quatro execuções, de modo que o ganho de macro-F1 provém
exclusivamente das outras duas classes. A segunda é que o filtro não pode elevar o
\emph{recall}, apenas a precisão, e a cobertura do léxico passa a ser um teto de
\emph{recall} da classe-alvo, somado ao teto já imposto pela janela
(Seção~\ref{sec:teto_recall}).

Convém precisar o que esta etapa é e o que ela não é. \textbf{Nenhum modelo é
retreinado aqui.} Os pesos empregados são os das execuções reportadas no
Capítulo~\ref{cap:experimentos}, e o RECLin-PT resultante \textbf{não possui pesos
próprios}. Trata-se de um sistema híbrido, em que um componente determinístico
restringe o espaço de decisão e o \emph{encoder} já ajustado escolhe, dentro desse
espaço, qual entidade é negada. O filtro tampouco representa um retorno a
abordagens puramente léxicas como o NegEx \cite{chapman_2001_negex}, pois ele não
decide a relação. Quem decide o alvo da negação continua sendo o classificador
contextual, e a Seção~\ref{sec:regra_pura} introduz o sistema de referência que
mede quanto a restrição léxica alcança sozinha.

\subsection{Calibração de \texttt{min\_freq} no conjunto de desenvolvimento}
\label{sec:calibracao}

O único parâmetro livre do filtro é o limiar \texttt{min\_freq}, e a sua escolha é
feita \textbf{no conjunto de desenvolvimento}, nunca no de teste. O critério foi
fixado antes da varredura. Seleciona-se o \texttt{min\_freq} que maximiza a
\emph{média} do F1 de \texttt{negation\_of} no desenvolvimento sobre as quatro
execuções, com empate resolvido pelo menor limiar, que corresponde ao léxico de
maior cobertura e, portanto, ao maior teto de \emph{recall}. A varredura sobre
$\{1,\,2,\,3,\,5,\,10\}$ seleciona \texttt{min\_freq}$=3$, com léxico de $11$
formas e cobertura de $0{,}9467$ dos pares \texttt{negation\_of} do
desenvolvimento.

Essa escolha de partição é o ponto metodologicamente decisivo da etapa, e o seu
custo é mensurável. Uma análise exploratória anterior, conduzida diretamente sobre
o conjunto de teste, encontrava ganho monotônico com léxicos cada vez mais
restritivos, e o seu melhor valor estava em \texttt{min\_freq}$=10$. No conjunto
de desenvolvimento, por sua vez, o máximo é \emph{interior} e recai sobre
\texttt{min\_freq}$=3$, porque a cobertura começa a diminuir antes que o ganho de
precisão a compense. A configuração metodologicamente defensável não é, portanto,
a que o conjunto de teste premiaria, e a distância entre as duas é exatamente o que
a disciplina de calibração custa. Registrar esse custo é parte do argumento, e não
um detalhe de implementação, uma vez que é ele que sustenta que o ganho reportado
no Capítulo~\ref{cap:experimentos} não decorre de ajuste ao conjunto de avaliação.
O conjunto de teste foi lido uma única vez por sistema, com a configuração já
fechada, conforme o protocolo da Seção~\ref{sec:avaliacao}.

A mesma calibração avaliou ainda um sistema combinado, em que a condição de
distância da regra descrita a seguir seria aplicada também sobre as predições já
filtradas. O limiar vencedor no conjunto de desenvolvimento desliga essa condição,
de modo que o sistema combinado colapsa no filtro isolado. Ele não foi, por esse
motivo, levado ao conjunto de teste como sistema separado.

\subsection{A regra pura como terceiro sistema de referência}
\label{sec:regra_pura}

A comparação entre dois \emph{encoders} não responde à pergunta que se formula em
seguida. O que o modelo contextual acrescenta sobre uma regra fixa? Para
respondê-la, define-se um terceiro sistema de referência que dispensa qualquer
rede neural e opera somente sobre o mesmo léxico induzido do treino
(\texttt{scripts/make\_rule\_baseline.py}). Quatro variantes foram consideradas,
todas partindo da condição de que $e_1$ seja um gatilho do léxico, e sendo
\texttt{gap} a distância em caracteres entre os \emph{spans}, a mesma grandeza
empregada na geração de candidatos:

\begin{enumerate}
  \item \textbf{R1}: todo par cujo $e_1$ é gatilho recebe \texttt{negation\_of}.
  \item \textbf{R2}: R1 restrita ao alvo mais próximo de cada gatilho, um alvo por
        gatilho.
  \item \textbf{R3}: R1 restrita aos pares cujo \texttt{gap} não excede um limiar.
  \item \textbf{R4}: R1 sob as duas restrições anteriores simultaneamente.
\end{enumerate}

A variante e o limiar também são escolhidos no conjunto de desenvolvimento, pelo
F1 de \texttt{negation\_of}, com empates resolvidos pela regra de menor índice e
depois pelo menor limiar. A escolha recai sobre \textbf{R3 com \texttt{gap}$\leq
1$}, que alcança F1 de $0{,}6935$ no desenvolvimento. A regra nunca prediz
\texttt{associated\_with} e não constitui, por isso, um sistema completo de
extração de relações. Ela entra como piso da classe-alvo, e o seu papel é delimitar
quanto do ganho do RECLin-PT seria atribuível à restrição lexical isolada.

As predições dos cinco sistemas introduzidos nesta seção são gravadas no
\textbf{espaço completo} de candidatos da partição, na mesma ordem e com o mesmo
\emph{gold} dos arquivos de predição das execuções dos \emph{baselines}. Essa
condição é o que permite compará-los sob o mesmo teste pareado aplicado no
Capítulo~\ref{cap:experimentos}, que aborta quando os rótulos de referência de
dois sistemas divergem.
